\scanpage{195}

Conclusion : les $u_{\deleted{\unsure{j}}i}$ \inserted{$(0\leq i\leq k-1)$} sont des bijections, les $E_i$ sont de dim. 1.

% slanted
La représentation de $\unsure{\mathfrak g|\mathscr G}$ ou de $G$ qu'on vient de construire est notée $e_{n,k+1}$ $(k+1\geq2)$ \inserted{\unsure{Or|Car} la repr. ainsi construite est indécomposable, car elle induit déjà sur $F$ la repr. indécomposable notée $e_{k+1}$ plus haut} Si $k=0$, il est naturel de poser $e_{n,1}=\sigma_n$.

Ainsi

\noindent\begin{tabular}{@{}|p{.92\linewidth}@{}}\uline{Proposition} Les représentations indécomposables de $G$ sont \deleted{les}\inserted{équivalentes aux} représentations $e_{n,k}$ \deleted{\unsure{$(k\geq1)$}} $(n\in Z,\ k\in Z,\ k\geq1)$ qu'on vient de construire : \deleted{\unsure{$\sigma$}} $e_{n,1}$ est la représentation $\sigma_n$ de dim 1 induite par la représentation de $=$ nom de $H=C^*$, $e_{n,k}=\sigma_n\otimes e_{0,k}$, où $e_{0\,k}$ est la représentation \deleted{\unsure{à val}} dans l'espace $k[Z]/Z^k$ définie ainsi : $X$ \unsure{définit|désigne|étant} l'opérateur de multiplication par $Z$ $\bmod (Z^k)$, et

\[Y.Z^\alpha=\alpha p Z^\alpha\quad(Z^k).\]

\inserted{les plus bas et plus hauts poids de $e_{n,k}$ sont $n$ et $n+p(k-1)$, donc} les représentations $e_{n,k}=\sigma_n\otimes e_{0\,k}$ sont deux à deux inéquivalentes. $e_{n,k}$ est de degré $k$ ; $\overset{\vee}{e}_{n,k}=e_{-n-p(k-1),k}$.\end{tabular}

% archive sheet number: 194
